\chapter{Study of Deformation}

If $R$ is the reference configuration of a body, let $\hat{\tenq{y}}(\bullet,t)\colon R \to R_t$. Choose any $t_d \in \tau$. Let $R_d\equiv R_{t_d} = \hat{\tenq{y}}(R,t_d)$, so we can write 
\[\hat{\tenq{y}}(\bullet,t_d) = \hat{\tenq{y}}(\bullet) \colon R \to R_d\]
where $\hat{\tenq{y}}(\bullet)$ is a deformation. \\

Define a \emph{displacement field} $\tenq{u}(\tenq{x})$ such that 
\[\hat{\tenq{y}}(\tenq{x}) = \tenq{x} + \tenq{u}(\tenq{x})\]

\begin{figure}[H]
\centering
\resizebox{0.4\textwidth}{!}{\input{figure8.eps_tex}}

\caption{Scheme of deformation field}
\label{Scheme of deformation field}
\end{figure}

\section{Deformation gradient field}

\[\tenq{F}(\tenq{x}) = \nabla \hat{\tenq{y}}(\tenq{x}) \in \mathscr{L}_{+}\]

where $\mathscr{L}_{+}$ is set of all $2^{nd}$ order tensor with positive determinant\\

or 

\[F_{ij} = \frac{\partial y_i(\tenq{x})}{\partial x_j} = y_{i,j} \text{ on }R\]
\[\det\tenq{F}(\tenq{x}) = \det \nabla \hat{\tenq{y}}(\tenq{x}) = J(\tenq{x})>0\]
\[\nabla \hat{\tenq{x}}(\tenq{y}) = \tenq{F}^{-1}(\hat{\tenq{x}}(\tenq{y}))\phantom{-}\forall\tenq{y}\in R_d\]
\[\because \hat{\tenq{y}}(\tenq{x}) = \tenq{x} + \tenq{u}(\tenq{x})\]
\[\Rightarrow \nabla \hat{\tenq{y}}(\tenq{x}) = \tenq{1} + \nabla \tenq{u}(\tenq{x})\]
\[\Rightarrow \tenq{F}(\tenq{x}) = \tenq{1} + \nabla \tenq{u}(\tenq{x})\]
or 
\[F_{ij}(\tenq{x}) = \delta_{ij} + u_{i,j}(\tenq{x})\]
where $\nabla \tenq{u}(\tenq{x})$ is called displacement gradient field.


\begin{lemma} \label{volume function for subregion}

Let $\hat{\tenq{y}}\colon R \to R_d$ be a deformation and let $\phi \colon R_d \to R$ be continuous. Then for any subregion (compact) $P_d \subset R_d$

\[\int_{P_d} \phi (\tenq{y}) d\volume_{\tenq{y}} =   \int_{P} \phi (\hat{\tenq{y}}(\tenq{x}))J(\tenq{x})d\volume_{\tenq{x}}\]

where 

\begin{align*}
P &= \hat{\tenq{y}}^{-1}(P_d), P_d = \hat{\tenq{y}}(P)\\
J &= \det \nabla \hat{\tenq{y}} \text{ on }R
\end{align*}
\end{lemma}

\begin{definition}

$\hat{\tenq{y}}\colon R \to R_d$ is \emph{locally volume preserving} (l.v.p.) or \emph{isochoric} if for \emph{every} compact $P \subset R$

\[\text{volume }P = \text{volume }P_d, P_d = \hat{\tenq{y}}(P)\]

\end{definition}


\begin{claim}

A definition $\hat{\tenq{y}}\colon R \to R_d$ is isochoric if and only if $J(\tenq{x}) = 1 \phantom{-}\forall \tenq{x} \in R$

\end{claim}

\begin{proof}\mbox{}\\

For any $P \subset R$ with $P_* = \hat{\tenq{y}}(P)\subset R_*$

\[\text{volume }(P) = \int_P d\volume_{\tenq{x}}\]
\[ \text{volume }(P_*) = \int_{P_*} d\volume_{\tenq{y}} \vc{=}{Lemma}{\eqref{volume function for subregion}} \int_P J(\tenq{x})d\volume_{\tenq{x}}\]

Hence $\hat{\tenq{y}}$ is l.v.p. if and only if 
\[\text{volume }(P) - \text{volume }(P_*) = 0\]
or 
\[\int_P \left[J(\tenq{x}) - 1 \right] d\volume_{\tenq{x}} = 0 \phantom{-}\forall P \subset R\]
\[\Rightarrow J(\tenq{x}) = 1\]

Define an incompressible body as that possible induced deformation are l.v.p.


\end{proof}

\begin{definition}{Rigid Deformation}

An admissible deformation $\hat{\tenq{y}}\colon R \to R_*$ is rigid if 

\[\abs{\hat{\tenq{y}}(\tenq{x}) - \hat{\tenq{y}}(\tenq{z})} = \abs{\tenq{x} - \tenq{z}}\phantom{-}\forall \tenq{x},\tenq{z} \in R\]

\begin{figure}[H]
\centering
\resizebox{0.7\textwidth}{!}{\input{figure9.eps_tex}}

\caption{Scheme of rigid deformation}
\label{Scheme of rigid deformation}
\end{figure}

\end{definition}

\begin{claim}

$\hat{\tenq{y}}$ is rigid if and only if 

\begin{equation}
\tenq{y} = \hat{\tenq{y}}(\tenq{x}) = \tenq{Q}\tenq{x} + \tenq{c} \tag{1} \label{rigid deformation}
\end{equation}

where \\

$\tenq{c}$ is a constant tensor \\
$\tenq{Q}$ is a constant proper orthogonal $2^{nd}$ order tensor, i.e. $\tenq{Q}\tenq{Q}^T = \tenq{1}, \det\tenq{Q} = +1$ 


\end{claim}

\begin{proof}\mbox{}\\

\q{$\Leftarrow$} Assume that \eqref{rigid deformation} holds. Then \eqref{rigid deformation} is admissible since


\begin{itemize}

\item 
\eqref{rigid deformation} is 1 to 1

\item 
\eqref{rigid deformation} is $C^{\infty}(R)$

\item From \eqref{rigid deformation}
\[\nabla \hat{\tenq{y}} = \tenq{F}(\tenq{x}) = \tenq{Q}\]
or 
\[y_{i,j} = \frac{\partial}{\partial x_j}\left(Q_{ik}x_k + c_i \right) = Q_{ij}\]
\[\Rightarrow \det \tenq{F} = \det \nabla \hat{\tenq{y}} = \det \tenq{Q} = +1 >0\]
\[\abs{\hat{\tenq{y}}(\tenq{x}) - \hat{\tenq{y}}(\tenq{z})} \vc{=}{\eqref{rigid deformation}}{} \abs{\tenq{Q}\left(\tenq{x} - \tenq{z}\right)} \vc{=}{}{\downarrow} \abs{\tenq{x} - \tenq{z}}\]

Recall $\tenq{u}\cdot\tenq{W}\tenq{v} = \left(\tenq{W}^T\tenq{u}\right)\cdot\tenq{v}$

\begin{align*}
\therefore\abs{\tenq{Q}\left(\tenq{x} - \tenq{z}\right)} 
&= \sqrt{\underbrace{\tenq{Q}\left(\tenq{x} - \tenq{z}\right)}_{\tenq{u}}\cdot\underbrace{\tenq{Q}}_{\tenq{W}}\underbrace{\left(\tenq{x} - \tenq{z}\right)}_{\tenq{v}}}\\
&= \sqrt{\underbrace{\tenq{Q}^T\tenq{Q}}_{\tenq{1}}\left(\tenq{x} - \tenq{z}\right)\cdot\left(\tenq{x} - \tenq{z}\right)}\\
&= \abs{\tenq{x} - \tenq{z}}
\end{align*}
\end{itemize}

\q{$\Rightarrow$} Assume that $\hat{\tenq{y}}$ is rigid then
\[\left[\hat{y}_i(\tenq{x}) - \hat{y}_i(\tenq{z}) \right]\left[\hat{y}_i(\tenq{x}) - \hat{y}_i(\tenq{z}) \right] - \left(x_i - z_i\right)\left(x_i - z_i\right) = 0\]

Take $\frac{\partial}{\partial x_j}$ of above
\[2F_{ij}(\tenq{x})\left[\hat{y}_i(\tenq{x}) - \hat{y}_i(\tenq{z}) \right] - 2\delta_{ij}\left(x_i - z_i\right) = 0\]
\[\Rightarrow\tenq{F}^T(\tenq{x})\left[\hat{\tenq{y}}(\tenq{x}) - \hat{\tenq{y}}(\tenq{z}) \right] = \tenq{x} - \tenq{z}\]
or 

\[2F_{ij}(\tenq{x})\left[\hat{y}_i(\tenq{x}) - \hat{y}_i(\tenq{z}) \right] - 2\delta_{ij}\left(x_i - z_i\right) = 0\]
\[\Rightarrow\tenq{F}^T(\tenq{x})\left[\hat{\tenq{y}}(\tenq{z}) - \hat{\tenq{y}}(\tenq{x}) \right] = \tenq{z} - \tenq{x}\]
Take gradient with respect to $\tenq{z}$

\begin{equation}
\tenq{F}^T(\tenq{x})\tenq{F}(\tenq{z}) = \tenq{1}\phantom{-}\forall\tenq{x},\tenq{z}\in R \tag{2} \label{F of rigid deformation}
\end{equation}

Set $\tenq{z} = \tenq{x}$

\[\tenq{F}^T(\tenq{x})\tenq{F}(\tenq{x}) = \tenq{1} \Rightarrow \tenq{F}(\tenq{x}) \text{ is orthogonal}\]

Fixed $\tenq{x}$ in \eqref{F of rigid deformation}, $\tenq{F}(\tenq{z})$ is constant tensor. \\
$\det \tenq{F}>0$ ($\because$ admissible) $\Rightarrow \tenq{F}$ is constant and proper orthogonal or $\tenq{F} = \tenq{Q}$. 

\end{proof}

\begin{note}
If $\tenq{v} = \tenq{x} - \tenq{z}$, $\tenq{v}_* = \hat{\tenq{y}}(\tenq{x}) - \hat{\tenq{y}}(\tenq{z})$, then 
$\tenq{v}_* = \tenq{Q}\tenq{v}$.
\end{note}

\begin{remark}\mbox{}\\

\begin{enumerate}[label=(\alph*)]
\item 
A reflection is distance preserving but not admissible ($\because \det \tenq{F}<0$) and thus not rigid. 

\item 
Rigid deformation is l.v.p. since $J = 1$.
\[J = \det \nabla \hat{\tenq{y}} = \det\tenq{F}(\tenq{x}) = \det \tenq{Q} = 1\]


\end{enumerate}
\end{remark}



\begin{definition}{Homogeneous Deformation}

A deformation $\hat{\tenq{y}}\colon R \to R_*$ is homogeneous if $\nabla \hat{\tenq{y}} = $ constant tensor on $R$, i.e. there exists $\tenq{F} = \text{ constant }\in \mathscr{L}_+$ (set of all $2^{nd}$ order tensors with positive determinant, $\det \tenq{F} >0$) and $\tenq{c}=\text{ constant }\in E_3$ such that

\[\tenq{y} = \hat{\tenq{y}} = \tenq{F}\tenq{x} + \tenq{c}\phantom{-}\forall\tenq{x} \in R\]

\end{definition}

Any general admissible deformation, $\tenq{y} = \hat{\tenq{y}}(\tenq{x})$, is homogeneous locally.\\
Let $\tenq{y} = \hat{\tenq{y}}(\tenq{x})$ be a general deformation (admissible)

\begin{figure}[H]
\centering
\resizebox{0.7\textwidth}{!}{\input{figure10.eps_tex}}

\caption{Scheme of homogeneous deformation locally}
\label{Scheme of homogeneous deformation locally}
\end{figure}

For any $\mathring{\tenq{x}} \in R$ by Taylor theorem and smoothness of any admissible deformation

\[\hat{\tenq{y}}(\tenq{x}) = \hat{\tenq{y}}(\mathring{\tenq{x}}) + \tenq{F}(\mathring{\tenq{x}})\left(\tenq{x} - \mathring{\tenq{x}}\right) + O\left(\abs{\tenq{x} - \mathring{\tenq{x}}}\right)\]

As $\abs{\tenq{x} - \mathring{\tenq{x}}}\to 0$
\[\hat{y}_i(\tenq{x}) = \hat{y}_i(\mathring{\tenq{x}}) + \frac{\partial \hat{y}_i(\mathring{\tenq{x}})}{\partial x_j}\left(x_j - \mathring{x}_j\right)\]
or 
\[\hat{\tenq{y}}(\tenq{x}) = \tenq{F}(\mathring{\tenq{x}})\tenq{x} + \hat{\tenq{y}}(\mathring{\tenq{x}}) - \tenq{F}(\mathring{\tenq{x}})\mathring{\tenq{x}}\]
i.e.

\[\hat{\tenq{y}}(\tenq{x}) = \tenq{F}\tenq{x} + \tenq{c}\]
where


\[\tenq{F} = \tenq{F}(\mathring{\tenq{x}}) = \text{ constant tensor}, \det\tenq{F}(\mathring{\tenq{x}})>0\]
\[\tenq{c} \buildrel \Delta \over =  \hat{\tenq{y}}(\mathring{\tenq{x}}) - \tenq{F}(\mathring{\tenq{x}})\mathring{\tenq{x}} = \text{ constant }\]

\begin{claim}\mbox{}\\

\begin{enumerate}[label=(\alph*)]

\item 
For a homogeneous deformation
\[\mathcal{H}\colon \tenq{y} = \hat{\tenq{y}}(\tenq{x}) = \tenq{F}\tenq{x} + \tenq{c}\]
$\det \tenq{F}>0, \tenq{x} \in R$, the inverse $\hat{\tenq{y}}^{-1}$ is also a homogeneous deformation. 

\item 
If further\footnotemark
\[\mathcal{H}'\colon \hat{\tenq{y}}'(\tenq{x}) = \tenq{F}'\tenq{x} + \tenq{c}',\det\tenq{F}'>0\]
\[\mathcal{H}''\colon \hat{\tenq{y}}''(\tenq{x}) = \tenq{F}''\tenq{x} + \tenq{c}'',\det\tenq{F}''>0\]
Then $\mathcal{H} = \mathcal{H}'\mathcal{H}''$ defined as 

\[\hat{\tenq{y}}(\tenq{x}) = \tenq{F}''(\hat{\tenq{y}}'(\tenq{x})) + \tenq{c}''\]
is also homogeneous. 

\footnotetext{$\mathcal{H}'$ and $\mathcal{H}''$ are not \emph{prime} , they mean the different deformation but in same motion.}

\begin{figure}[H]
\centering
\resizebox{0.7\textwidth}{!}{\input{figure11.eps_tex}}

\caption{Scheme of homogeneous deformation continuously}
\label{Scheme of homogeneous deformation continuously}
\end{figure}


\end{enumerate}

\end{claim}

\begin{proof}\mbox{}\\

\begin{enumerate}[label=(\alph*)]

\item 
\[\hat{\tenq{y}} = \tenq{F}\tenq{x} + \tenq{c}\]
\[\Rightarrow\tenq{F}^{-1}\hat{\tenq{y}} = \tenq{x} + \tenq{F}^{-1}\tenq{c}\]
\[\therefore \tenq{x} = \hat{\tenq{x}}(\tenq{y}) = \tenq{F}^*\tenq{y} + \tenq{c}^*\]

where 

\begin{align*}
\tenq{F}^* 
&= \tenq{F}^{-1} = \text{ constant and }\det\tenq{F}^* = \det\tenq{F}^{-1} > 0\footnotemark\\
\tenq{c}^* &= -\tenq{F}^{-1}\tenq{c} = \text{ constant}
\end{align*}
\footnotetext{$\because \det \tenq{F}>0$}
\item 
\[\hat{\tenq{y}}(\tenq{x}) = \tenq{F}''\left(\tenq{F}'\tenq{x} + \tenq{c}'\right) + \tenq{c}'' = \tenq{F}\tenq{x} + \tenq{c}\]

where 

\begin{align*}
\tenq{F} &= \tenq{F}''\tenq{F}' = \text{ constant},\det\tenq{F} = \det\tenq{F}''\det\tenq{F}'>0\\
\tenq{c} &= \tenq{F}''\tenq{c}' + \tenq{c}'' = \text{ constant}
\end{align*}

\end{enumerate}
\end{proof}

\begin{figure}[H]
\centering
\resizebox{0.7\textwidth}{!}{\input{figure12.eps_tex}}

\caption{Scheme of homogeneous deformation}
\label{Scheme of homogeneous deformation}
\end{figure}

\begin{align*}
\tenq{v} &= \tenq{z} - \tenq{x}\\
\tenq{v}^* 
& = \hat{\tenq{y}}(\tenq{z}) - \hat{\tenq{y}}(\tenq{x})\\
& = \left(\tenq{F}\tenq{z} + \tenq{c}\right) - \left(\tenq{F}\tenq{x} + \tenq{c}\right)\\
& = \tenq{F}\left(\tenq{z}-\tenq{x}\right) = \tenq{F}\tenq{v}
\end{align*}

\begin{equation}
\boxed{\tenq{v}^* = \tenq{F}\tenq{v} \text{ or }v_i^* = F_{ij}v_j}\tag{1} \label{homogeneous deformation}
\end{equation} 

For a rigid deformation, $\tenq{F} = \tenq{Q}$, $v_i^* = Q_{ij}v_j$, i.e. pure rotation. 

\begin{theorem}\mbox{}\\

Let $\tenq{v}$, $\tenq{w}$ be vectors in $R$, $\tenq{v}^*$, $\tenq{w}^*$ be their images under homogeneous deformation

\[\mathcal{H}\colon \hat{\tenq{y}}(\tenq{x}) = \tenq{F}\tenq{x} + \tenq{c}\]
Then 

\begin{equation}
\tenq{v}^*\cdot\tenq{w}^* - \tenq{v}\cdot\tenq{w} = 2\tenq{v}\cdot\left(\tenq{E}\tenq{w}\right) = 2v_iE_{ij}w_j = 2E_{ij}v_iw_j \tag{2} \label{homogeneous deformation- dot}
\end{equation}

where 

\begin{equation}
\tenq{E} = \frac{1}{2}\left\{ \tenq{F}^T\tenq{F}  - \tenq{1}\right\} \cdots\text{Lagrangian (Green) strain tensor} \tag{3} \label{Lagrangian (Green) strain tensor}
\end{equation}

or 

\begin{align*}
E_{ij} &= \frac{1}{2}\left\{F_{ik}^TF_{kj} - \delta_{ij} \right\} = \frac{1}{2}\left\{F_{ki}F_{kj} - \delta_{ij} \right\} \\
&= \frac{1}{2}\left\{ \frac{\partial y_k}{\partial x_i}\frac{\partial y_k}{\partial x_j} - \delta_{ij}\right\}
\end{align*}
\[\cdots \text{Components of Lagrangian (Green) strain tensor}\]

\begin{proof}

\begin{align*}
\tenq{v}^*\cdot\tenq{w}^* - \tenq{v}\cdot\tenq{w} 
&\vc{=}{\eqref{homogeneous deformation}}{}\tenq{F}\tenq{v}\cdot\tenq{F}\tenq{w} - \tenq{v}\cdot\tenq{w}\\
& = \footnotemark \tenq{F}^T\tenq{F}\tenq{w}\cdot\tenq{v} - \tenq{v}\cdot\tenq{w}\\
& = \tenq{v}\cdot\tenq{F}^T\tenq{F}\tenq{w} - \tenq{v}\cdot\tenq{w}\\
& = \tenq{v}\cdot\left(\tenq{F}^T\tenq{F} - \tenq{1}\right)\tenq{w}\\
&\vc{=}{\eqref{homogeneous deformation- dot}}{}2\tenq{v}\cdot\tenq{E}\tenq{w}
\end{align*}
\footnotetext{$\tenq{F}\tenq{w}\cdot\tenq{F}\tenq{v} = \tenq{F}^T\tenq{F}\tenq{w}\cdot\tenq{v}$}

\end{proof}

When $\tenq{F} = \tenq{Q}$ (rigid deformation) $\Rightarrow\tenq{E} = \tenq{0}$

\begin{note}

\begin{align*}
\tenq{v}^*\cdot\tenq{w}^* - \tenq{v}\cdot\tenq{w} 
&\vc{=}{\eqref{homogeneous deformation}}{}\tenq{v}^*\cdot\tenq{w}^* - \tenq{F}^{-1}\tenq{v}^*\cdot\tenq{F}^{-1}\tenq{w}^*\\
& = \tenq{v}^*\cdot\tenq{w}^* - \tenq{v}^*\cdot\tenq{F}^{-T}\tenq{F}^{-1}\tenq{w}^*\\
& = \tenq{v}^*\cdot\left(\tenq{1}-\tenq{F}^{-T}\tenq{F}^{-1}\right)\tenq{w}\\
&\vc{=}{\Delta}{}2\tenq{v}^*\cdot\tenq{M}\tenq{w}^*
\end{align*}
where 
\[\tenq{M} = \frac{1}{2}\left\{\tenq{1} - \tenq{F}^{-T}\tenq{F}^{-1} \right\} \cdots\text{Eulerian (Almansi) strain tensor}\]
or 

\begin{align*}
M_{ij} &= \frac{1}{2}\left\{\delta_{ij} - F_{ik}^{-T}F_{kj}^{-1} \right\} \vc{=}{Hwk}{} \frac{1}{2}\left\{\delta_{ij} - F_{ki}^{-1}F_{kj}^{-1} \right\} \\
&= \frac{1}{2}\left\{ \delta_{ij} - \frac{\partial x_k}{\partial y_i}\frac{\partial x_k}{\partial y_j} \right\}
\end{align*}
\[\cdots \text{Components of Eulerian (Almansi) strain tensor}\]

\begin{itemize}

\item 
Both $\tenq{E}$ and $\tenq{M}$ are symmetric. 

\item 
For rigid deformation, $\tenq{F} = \tenq{Q}\Rightarrow \tenq{E} = \tenq{M} = \tenq{0}$

\end{itemize}

\end{note}



\end{theorem}



\section{Strain tensor in terms of displacement field}

Recall that 

\begin{align*}
\hat{\tenq{y}}(\tenq{x}) &= \tenq{x} + \tenq{u}(\tenq{x})\Rightarrow \tenq{F} = \nabla_{\tenq{x}}\hat{\tenq{y}}(\tenq{x}) = \tenq{1} + \nabla_{\tenq{x}}\tenq{u}(\tenq{x})\\
\hat{\tenq{x}}(\tenq{y}) &= \tenq{y} - \tenq{u}(\tenq{y})\Rightarrow \tenq{F}^{-1} = \nabla_{\tenq{y}}\hat{\tenq{x}}(\tenq{y}) = \tenq{1} - \nabla_{\tenq{y}}\tenq{u}(\tenq{y})\\
\end{align*}

\begin{align*}
\tenq{E} 
& = \frac{1}{2}\left\{ \tenq{F}^T\tenq{F}  - \tenq{1}\right\} \\
& = \frac{1}{2}\left\{\left[\tenq{1} + \nabla_{\tenq{x}}\tenq{u}(\tenq{x})\right]^T\left[\tenq{1} + \nabla_{\tenq{x}}\tenq{u}(\tenq{x})\right] - \tenq{1}\right\}\\
& = \frac{1}{2}\left\{\nabla_{\tenq{x}}\tenq{u} + \left(\nabla_{\tenq{x}}\tenq{u}\right)^T + \left(\nabla_{\tenq{x}}\tenq{u}\right)^T\left(\nabla_{\tenq{x}}\tenq{u}\right)\right\}
\end{align*}

or 

\[E_{ij} = \frac{1}{2}\left(\frac{\partial u_i}{\partial x_j} + \frac{\partial u_j}{\partial x_i} + \frac{\partial u_k}{\partial x_i}\frac{\partial u_k}{\partial x_j}\right)\]

\begin{align*}
\tenq{M} 
& = \frac{1}{2}\left\{\tenq{1} -  \tenq{F}^{-T}\tenq{F}^{-1}\right\} \\
& = \frac{1}{2}\left\{\tenq{1} - \left[\tenq{1} - \nabla_{\tenq{y}}\tenq{u}(\tenq{y})\right]^T\left[\tenq{1} - \nabla_{\tenq{y}}\tenq{u}(\tenq{y})\right]\right\}\\
& = \frac{1}{2}\left\{\nabla_{\tenq{y}}\tenq{u} + \left(\nabla_{\tenq{y}}\tenq{u}\right)^T - \left(\nabla_{\tenq{y}}\tenq{u}\right)^T\left(\nabla_{\tenq{y}}\tenq{u}\right)\right\}
\end{align*}

or 

\[M_{ij} = \frac{1}{2}\left(\frac{\partial u_i}{\partial y_j} + \frac{\partial u_j}{\partial y_i} - \frac{\partial u_k}{\partial y_i}\frac{\partial u_k}{\partial y_j}\right)\]

From polar decomposition theorem

\begin{align*}
\tenq{F} & = \tenq{Q}\tenq{U}\cdots\text{right decomposition}\\
& = \tenq{V}\tenq{Q}\cdots\text{left decomposition}
\end{align*}
with 
\[\tenq{U} = \sqrt{\tenq{F}^T\tenq{F}}, \tenq{V} = \sqrt{\tenq{F}\tenq{F}^T}, \tenq{Q} = \tenq{F}\tenq{U}^{-1}= \tenq{V}^{-1}\tenq{F}\]

which is proper orthogonal, $\because \det\tenq{F}>0$ for admissible deformation. 

\begin{definition}\mbox{}\\

\begin{align*}
\tenq{C} &\vc{=}{\Delta}{}\tenq{F}^T\tenq{F}\text{ called right deformation tensor}\\
\tenq{G} &\vc{=}{\Delta}{}\tenq{F}\tenq{F}^T\text{ called left deformation tensor}
\end{align*}
Then 
\[\tenq{C} = \tenq{U}^2, \tenq{G} = \tenq{V}^2, \tenq{E} = \frac{1}{2}\left\{\tenq{C}-\tenq{1}\right\}\]

According to the polar decomposition theorem, $\tenq{C}$ and $\tenq{G}$ are both symmetric and positive definite and have the same eigenvalues but different eigenvectors, i.e. different principle frame. 

\end{definition}

\begin{proof}\mbox{}\\
Let $X = \left\{ {0;\tenq{e}_1,\tenq{e}_2,\tenq{e}_3} \right\} \in \mathscr{F}$ be the principal frame for $\tenq{G}$ and let $\lambda_i>0$ be the three eigenvalues

\[\tenq{G}\tenq{e}_i = \lambda_i\tenq{e}_i\text{ (no sum)}\]

Let $X' = \left\{ {0;\tenq{e}_1',\tenq{e}_2',\tenq{e}_3'} \right\} \in \mathscr{F}$ with $\tenq{e}_i' = \tenq{Q}^T\tenq{e}_i$ such that

\[\tenq{C} = \tenq{Q}^T\tenq{G}\tenq{Q}\]

then 

\[\tenq{C}\tenq{e}_i' = \tenq{Q}^T\tenq{G}\tenq{Q}\tenq{Q}^T\tenq{e}_i = \tenq{Q}^T\tenq{G}\tenq{e}_i = \lambda_i\tenq{Q}^T\tenq{e}_i = \lambda_i\tenq{e}_i'\]

\[\Rightarrow\tenq{C}\tenq{e}_i' = \lambda_i\tenq{e}_i'\]

\end{proof}


From \eqref{homogeneous deformation- dot} and \eqref{Lagrangian (Green) strain tensor}

\begin{align*}
\tenq{v}^*\cdot\tenq{w}^* 
& = \tenq{v}\cdot\tenq{w} + 2\tenq{v}\cdot\tenq{E}\tenq{w}\\
& = \tenq{v}\cdot\tenq{w} + \tenq{v}\cdot\tenq{F}^T\tenq{F}\tenq{w} - \tenq{v}\cdot\tenq{w}
\end{align*}

\begin{equation}
\Rightarrow \boxed{\tenq{v}^*\cdot\tenq{w}^* = \tenq{v}\cdot\tenq{C}\tenq{w}} \tag{4} \label{deformation- dot}
\end{equation}

For rigid deformation, $\tenq{C} = \tenq{G} = \tenq{1}$


\section{Geometric interpretation of \texorpdfstring{$\tenq{E}$ or $E_{ij}^X$}{Lg}}


Define 

\begin{align*}
\lambda(\tenq{v}) &\vc{=}{\Delta}{}\frac{\abs{\tenq{v}^*}}{\abs{\tenq{v}}}\cdots\text{stretch ratio of }\tenq{v}\\
\epsilon(\tenq{v}) &\vc{=}{\Delta}{}\frac{\abs{\tenq{v}^*} - \abs{\tenq{v}}}{\abs{\tenq{v}}}\cdots\text{extension ratio of }\tenq{v}
\end{align*}

\[\Rightarrow \epsilon(\tenq{v}) = \lambda(\tenq{v}) - 1\]

If the deformation is homogeneous, then 

\[\lambda(\alpha\tenq{v}) = \lambda(\tenq{v}), \epsilon(\alpha\tenq{v}) = \epsilon(\tenq{v}), \alpha\in\mathbb{R}, \alpha\neq0\]

Set $\tenq{w} = \tenq{v}$ in \eqref{homogeneous deformation- dot}
\begin{align*}
\tenq{v}^*\cdot\tenq{v}^* - \tenq{v}\cdot\tenq{v} 
& = 2\tenq{v}\cdot\tenq{E}\tenq{v}\\
\Rightarrow \abs{\tenq{v}^*}^2 - \abs{\tenq{v}}^2 
& = 2\tenq{v}\cdot\tenq{E}\tenq{v}\\
\Rightarrow \frac{\abs{\tenq{v}^*}^2 - \abs{\tenq{v}}^2}{\abs{\tenq{v}}^2}
& = \frac{2\tenq{v}\cdot\tenq{E}\tenq{v}}{\abs{\tenq{v}}^2}\\
\Rightarrow \lambda^2(\tenq{v}) - 1 
& = \frac{2\tenq{v}\cdot\tenq{E}\tenq{v}}{\abs{\tenq{v}}^2}
\end{align*}

\[\Rightarrow \lambda(\tenq{v}) = \sqrt{1 + \frac{2\tenq{v}\cdot\tenq{E}\tenq{v}}{\abs{\tenq{v}}^2}} \phantom{-} \forall \tenq{v}\neq 0\]

Let $X = \left\{ {0;\tenq{e}_1,\tenq{e}_2,\tenq{e}_3} \right\} \in \mathscr{F}$

\begin{align*}
\epsilon(\tenq{e}_i) &= \lambda(\tenq{e}_i)-1\\
&= \sqrt{1 + \frac{2\tenq{e}_i\cdot\tenq{E}\tenq{e}_i}{1}} - 1 \text{ (no sum)}\\
&= \sqrt{1 + 2E_{ii}} - 1 \text{ (no sum)}
\end{align*}

\[E_{ii}^X = \frac{\lambda^2(\tenq{e}_i) - 1}{2}\text{ (no sum)}\]

So given a frame $X = \left\{ {0;\tenq{e}_1,\tenq{e}_2,\tenq{e}_3} \right\} \in \mathscr{F}$, the diagonal elements $E_{11}^X, E_{22}^X, E_{33}^X$ of the Lagrangian strain tensor in the frame are associated with the length changes of the unit vector of $X$.\\

Interpretation of $E_{ij}^X$, $i\neq j$

\begin{figure}[H]
\centering
\resizebox{0.7\textwidth}{!}{\input{figure13.eps_tex}}

\caption{Scheme of change of theta}
\label{Scheme of change of theta}
\end{figure}

\begin{align*}
\tenq{v}^*\cdot\tenq{w}^* &= \abs{\tenq{v}^*}\abs{\tenq{w}^*}\cos(\theta^*(\tenq{v}, \tenq{w}))\\
& \vc{=}{\eqref{homogeneous deformation- dot}}{} \tenq{v}\cdot\tenq{w} + 2\tenq{v}\cdot\tenq{E}\tenq{w}
\end{align*}
Also

\begin{align*}
\abs{\tenq{v}^*} &= \abs{\tenq{v}}\lambda(\tenq{v}) = \abs{\tenq{v}}\left[1 + \epsilon (\tenq{v}) \right]\\
\abs{\tenq{w}^*} &= \abs{\tenq{w}}\lambda(\tenq{w}) = \abs{\tenq{w}}\left[1 + \epsilon (\tenq{w}) \right]
\end{align*}

\begin{align}
\therefore \cos(\theta^*(\tenq{v}, \tenq{w})) 
&= \frac{\tenq{v}\cdot\tenq{w} + 2\tenq{v}\cdot\tenq{E}\tenq{w}}{\abs{\tenq{v}^*}\abs{\tenq{w}^*}}\nonumber \\
&= \frac{\tenq{v}\cdot\tenq{w} + 2\tenq{v}\cdot\tenq{E}\tenq{w}}{\abs{\tenq{v}}\abs{\tenq{w}}\left[1 + \epsilon (\tenq{v}) \right]\left[1 + \epsilon (\tenq{w}) \right]}\tag{$\ast$} \label{cos of deformation} 
\end{align}


Take $\tenq{v} = \tenq{e}_i$, $\tenq{w} = \tenq{e}_j$, $i\neq j$
\begin{figure}[H]
\centering
\resizebox{0.7\textwidth}{!}{\input{figure14.eps_tex}}

\caption{Deformation under principal frame}
\label{Deformation under principal frame}
\end{figure}

\[\therefore \cos(\theta^*(\tenq{e}_i, \tenq{e}_j)) = \frac{2\tenq{e}_i\cdot\tenq{E}\tenq{e}_j}{\sqrt{1+2E_{ii}}\sqrt{1+2E_{jj}}}, i \neq j\text{ (no sum)}\]

Define 
\[\mu_{ij}\vc{=}{\Delta}{}\frac{\pi}{2} - \theta^*(\tenq{e}_i, \tenq{e}_j)\]
Then

\[
\boxed{\sin\mu_{ij} = \frac{2E_{ij}^X}{\sqrt{1+2E_{ii}}\sqrt{1+2E_{jj}}}, i \neq j\text{ (no sum)}}
\]

\[\therefore \text{If } E_{ij}^X = 0,\forall i\neq j\]
\[\mu_{ij} = 0\text{ or }\theta^*(\tenq{e}_i, \tenq{e}_j) = \frac{\pi}{2}\]
i.e. 
\[\tenq{e}_i^*\perp\tenq{e}_j^*\]

Let $X \in \mathscr{F}$ be principal frame for $\tenq{C}$ and therefore principal frame for $\tenq{E}$. ($\because \tenq{C} = 2\tenq{E} + \tenq{1}$) Let $\lambda_i^2$ be eigenvalue for $\tenq{C}$. Then 

\[
[C_{ij}^X]
= 
\begin{bmatrix}
\lambda_1^2 & 0 & 0\\
0 & \lambda_2^2 & 0\\
0 & 0 & \lambda_3^2
\end{bmatrix}
=
\begin{bmatrix}
1+2E_1 & 0 & 0\\
0 & 1+2E_2 & 0\\
0 & 0 & 1+2E_3
\end{bmatrix}
\tag{$\ast$} \label{C of deformation}
\]
where $E_i$ are eigenvalue of $\tenq{E}$\\
What's the physical meaning of $\lambda_i$?\\
From $\eqref{C of deformation}$

\begin{align*}
\underbrace{\lambda_i}_{\text{mathematical stuff}} 
& = \sqrt{1 + 2E_i}\vc{=}{X\colon principal frame}{}
\sqrt{1+2E_{ii}^X}\\
& = \underbrace{\lambda(\tenq{e}_i^X)}_{\text{physical meaning}}
\end{align*}

\begin{figure}[H]
\centering
\resizebox{0.7\textwidth}{!}{\input{figure15.eps_tex}}

\caption{Deformation under principal frame}
\label{Deformation under principal frame-2}
\end{figure}

\begin{note}

The cube and the parallelepiped may have different orientation. 

\end{note}

\section{Volume change}

\[\volume(R_\ast) = \int_{R_\ast} d\volume_{\tenq{y}} = \int_{R}J d\volume_{\tenq{x}}  \vc{=}{homogeneous
}{deformation}J \int_{R}d\volume_{\tenq{x}} = Jd\volume(R)\]
\[J = \frac{\volume(R_\ast)}{\volume(R)}\]
From previous schematic Figure. \ref{Deformation under principal frame-2}, $J = \lambda_1\lambda_2\lambda_3$\\

Define \emph{volume dilatation} as 

\begin{align*}
\Theta 
&= \frac{\volume(R_\ast) - \volume(R)}{\volume(R)}\\
&= \lambda_1\lambda_2\lambda_3 - 1 = \sqrt{\left(1 + 2E_1\right)\left(1 + 2E_2\right)\left(1 + 2E_3\right)} - 1\\
&= \footnotemark \sqrt{1+2I_1(\tenq{E}) + 4I_2(\tenq{E}) + 8I_3(\tenq{E})} - 1
\end{align*}

\footnotetext{$I_1(\tenq{E}) = E_1 + E_2 + E_3, I_2(\tenq{E}) = E_1E_2 + E_2E_3 + E_1E_3, I_3(\tenq{E}) = E_1E_2E_3$}

\begin{definition}

$\mathcal{H}$ is a \emph{pure} homogeneous deformation if it is of the form

\[\mathcal{H} \colon \tenq{y} = \hat{\tenq{y}}(\tenq{x}) = \tenq{F}\tenq{x},\tenq{F} = \tenq{F}^T \text{ and }\tenq{F} \text{ is positive definite}\]
\end{definition}

\paragraph{Properties of the pure homogeneous definition}\mbox{}\\
In general\\
\[
\left.\begin{aligned}
\tenq{E} &= \frac{1}{2}\left(\tenq{F}^T\tenq{F} - \tenq{1}\right)\\
\tenq{C} &= \tenq{F}^T\tenq{F}\\
\tenq{G} &= \tenq{F}\tenq{F}^T
\end{aligned}\right\}
\Rightarrow
\begin{aligned}
\tenq{E} = \frac{1}{2}\left(\tenq{F}^2 - \tenq{1}\right)\\
\left.
\begin{aligned}
\tenq{C} &= \tenq{F}^2\\
\tenq{G} &= \tenq{F}^2
\end{aligned}
\right\}\Rightarrow
\tenq{C} = \tenq{G}
\end{aligned}
\]
i.e. 
\[\tenq{C},\tenq{G},\tenq{E},\tenq{U},\tenq{V} \text{ have same principal axes} \]

\[\tenq{C}, \tenq{G} \cdots \text{same principal values}\]
\[\tenq{U}, \tenq{V} \cdots \text{same principal values}\]

Consider $X \in \mathscr{F}$ (principal frame for $\tenq{C},\tenq{G},\tenq{E}$)


\[
[C_{ij}^X]
= 
[G_{ij}^X]
=
\begin{bmatrix}
\lambda_1^2 & 0 & 0\\
0 & \lambda_2^2 & 0\\
0 & 0 & \lambda_3^2
\end{bmatrix}
,
[F_{ij}^X]
=
\begin{bmatrix}
\lambda_1 & 0 & 0\\
0 & \lambda_2 & 0\\
0 & 0 & \lambda_3
\end{bmatrix}
\]

\begin{align*}
\mathcal{H} \colon \tenq{y} = \tenq{F}\tenq{x}
& \Rightarrow y_i^X = F_{ij}^X x_j^X\\
& \Rightarrow y_i^X = \lambda_{i} x_i^X\text{ (no sum)}
\end{align*}

\begin{figure}[H]
\centering
\resizebox{0.7\textwidth}{!}{\input{figure16.eps_tex}}

\caption{Scheme of pure homogeneous deformation}
\label{Scheme of pure homogeneous deformation}
\end{figure}

\begin{note}
The cube and the rectangular parallelepiped have the \emph{same} orientation. 
\end{note}

\begin{theorem}{Resolutions of Arbitrary Homogeneous Deformation}\mbox{}\\
Let $\tenq{y} = \hat{\tenq{y}}(\tenq{x}) = \tenq{F}\tenq{x} + \tenq{c}$, $\det\tenq{F}>0$\\
Recall $\tenq{F} = \tenq{Q}\tenq{U} = \tenq{V}\tenq{Q}$\\
where $\tenq{U},\tenq{V}$ are symmetric positive definite, $\tenq{Q}$ is orthogonal.

\begin{alignat*}{4}
& \mathcal{H}_1' \colon    \quad && \tenq{x}\to\tenq{y}'\colon 	\quad && \tenq{y}' = \tenq{U}\tenq{x}	    \quad && \text{(pure homogeneous } \\
& \phantom{-}  \quad && \phantom{-}	\quad && \phantom{-}	\quad && \text{deformation)} \\
& \mathcal{H}_1'' \colon   \quad && \tenq{y}'\to\tenq{y}''\colon 	\quad && \tenq{y}'' = \tenq{Q}\tenq{y}'	    \quad && \text{(rigid deformation about O)}\\
& \phantom{-}  \quad && \phantom{-}	\quad && \phantom{-}	\quad && \text{(rotation)} \\	
& \mathcal{H}_1''' \colon  \quad && \tenq{y}''\to\tenq{y}\colon	\quad && \tenq{y} = \tenq{y}'' + \tenq{c}	\quad && \text{(translation)} 		
\end{alignat*}

Then $\mathcal{H}\colon \tenq{x}\to \tenq{y}\colon\tenq{y} = \tenq{F}\tenq{x}+\tenq{c}$ is given by 

\begin{equation}
\mathcal{H} = \mathcal{H}_1'''\mathcal{H}_1''\mathcal{H}_1' \tag{$\ast$} \label{arbitary homogeneous deformation_1}
\end{equation}

Also suppose

\begin{alignat*}{4}
& \mathcal{H}_2' \colon    \quad && \tenq{x}\to\tenq{y}'\colon 	\quad && \tenq{y}' = \tenq{Q}\tenq{x}	    \quad && \text{(rigid deformation about O)}\\
& \phantom{-}  \quad && \phantom{-}	\quad && \phantom{-}	\quad && \text{(rotation)} \\
& \mathcal{H}_2'' \colon   \quad && \tenq{y}'\to\tenq{y}''\colon 	\quad && \tenq{y}'' = \tenq{V}\tenq{y}'	    \quad && \text{(pure homogeneous } \\
& \phantom{-}  \quad && \phantom{-}	\quad && \phantom{-}	\quad && \text{deformation)} \\
& \mathcal{H}_2''' \colon  \quad && \tenq{y}''\to\tenq{y}\colon	\quad && \tenq{y} = \tenq{y}'' + \tenq{c}	\quad && \text{(translation)} 		
\end{alignat*}

Then $\mathcal{H}\colon \tenq{x}\to \tenq{y}\colon\tenq{y} = \tenq{F}\tenq{x}+\tenq{c}$ is given by 

\begin{equation}
\mathcal{H} = \mathcal{H}_2'''\mathcal{H}_2''\mathcal{H}_2' \tag{$\ast\ast$} \label{arbitary homogeneous deformation_2}
\end{equation}
\end{theorem}

\begin{proof}\mbox{}\\
\[\tenq{y}' = \tenq{U}\tenq{x}, \tenq{y}'' = \tenq{Q}\tenq{y}'\]
\begin{align*}
\tenq{y} 
&= \tenq{y}''+\tenq{c}\\
&= \tenq{Q}\tenq{U}\tenq{x} + \tenq{c}\\
&= \tenq{F}\tenq{x} + \tenq{c} \Rightarrow \eqref{arbitary homogeneous deformation_1} \text{ holds}
\end{align*}
similar for $\eqref{arbitary homogeneous deformation_2}$. 
\end{proof}


\paragraph{2-D illustration of the two Equivalent Resolution of $\mathcal{H}$}\mbox{}\\
Let $X\in\mathscr{F}$ be principal frame for $\tenq{U},\tenq{C},\tenq{Q}$ and $X'\in\mathscr{F}$ be principal frame for $\tenq{V},\tenq{G}$.

\begin{figure}[H]
\centering
\resizebox{0.9\textwidth}{!}{\input{figure17.eps_tex}}

\caption{Scheme of Resolution of $\mathcal{H}_1$}
\label{Scheme of Resolution of $H_1$}
\end{figure}
\begin{figure}[H]
\centering
\resizebox{0.9\textwidth}{!}{\input{figure18.eps_tex}}

\caption{Scheme of Resolution of $\mathcal{H}_2$}
\label{Scheme of Resolution of $H_2$}
\end{figure}

Recall that $\tenq{V} = \tenq{Q}\tenq{U}\tenq{Q}^T$, $\tenq{U} = \tenq{Q}^T\tenq{U}\tenq{Q}$. If $\tenq{U}$ is symmetric and $\tenq{Q}$ is proper orthogonal and $X = \left\{ {0;\tenq{e}_1,\tenq{e}_2,\tenq{e}_3} \right\} \in \mathscr{F}$ is principal frame of $\tenq{U}$, $X' = \left\{ {0;\tenq{e}_1',\tenq{e}_2',\tenq{e}_3'} \right\} \in \mathscr{F}$ where $\tenq{e}_1' = \tenq{Q}\tenq{e}_1$ is principal frame for $\tenq{V} = \tenq{Q}\tenq{U}\tenq{Q}^T$.

\section{Direct geometrical interpretation of the Lagrangian Strain tensor field of an admissible deformation (General Deformation)}

Let $\hat{\tenq{y}}\colon R \to R_d$ be an admissible deformation, $\mathring{\tenq{x}}$ an interior point of $R$, and $\mathring{\tenq{\phi}} = \tenq{\phi}(\mathring{\tenq{x}})$ the Lagrangian stain tensor of the deformation.\\
Let further $X = \left\{ {0;\tenq{e}_1,\tenq{e}_2,\tenq{e}_3} \right\} \in \mathscr{F}$. 

\begin{enumerate}[label=(\alph*)]
\item 
Suppose $\Gamma_i$ ($i$ fixed) is a regular arc issuing from $\mathring{\tenq{x}}$, let $s$ be the arc-length along $\Gamma_i$ measured from $\mathring{\tenq{x}}$ and let $\tenq{e}_i$ be tangent to the $\Gamma_i$ at $\mathring{\tenq{x}}$. \\

Thus assume: \\
Parametrise $\Gamma_i$

\begin{equation}
\begin{split}
\Gamma_i \colon \tenq{x} &= \bar{\tenq{x}}(s)\phantom{-}(0\leq s\leq l), \bar{\tenq{x}}\in \footnotemark C^1([0,l])
,\\ 
\bar{\tenq{x}}(0)&= \mathring{\tenq{x}}, \bar{\tenq{x}}'(0)=\tenq{e}_i
\end{split}\tag{1} \label{Parametrise_1}
\end{equation}

\footnotetext{unique tangent exists}

Let $\accentset{\ast}{\Gamma}_i = \hat{\tenq{y}}(\Gamma_i)\cdots$deformation image of $\Gamma_i$ and $\accentset{\ast}{s}(s)$ the arc-length along $\accentset{\ast}{\Gamma}_i$ measured from $\mathring{\tenq{y}} = \hat{\tenq{y}}(\mathring{\tenq{x}})$.
Thus

\begin{equation}
\accentset{\ast}{\Gamma}_i \colon \tenq{y}=\hat{\tenq{y}}(\bar{\tenq{x}}(s)) = \bar{\tenq{y}}(s) \phantom{-}(0\leq s\leq l)\tag{2} \label{Parametrise_2}
\end{equation}

Then the arc-length along $\accentset{\ast}{\Gamma}_i$ is 

\begin{equation}
\accentset{\ast}{s}(s) = \int_0^s\left\{\bar{\tenq{y}}'(\sigma)\cdot \bar{\tenq{y}}'(\sigma) \right\}^{1/2}d\sigma \tag{3} \label{Parametrise_3}
\end{equation}
and 
\[\accentset{\ast}{s}'(0)\equiv {\left. {\frac{{d\accentset{\ast}{s}}}{{ds}}} \right|_{s = 0}} = \sqrt{1+2\phi_{ii}(\mathring{\tenq{x}})} \text{ (no sum)}\]

Note also that ${\left. {\frac{{d\accentset{\ast}{s}}}{{ds}}} \right|_{s = 0}} = \lambda(\mathring{\tenq{x}},\tenq{e}_i)$, where $\lambda(\mathring{\tenq{x}},\tenq{e}_i)$ is the stretch ratio of $\tenq{e}_i$ in the approximating local homogeneous deformation of $\mathring{\tenq{x}}$.

\begin{proof}\mbox{}\\

\begin{figure}[H]
\centering
\resizebox{0.9\textwidth}{!}{\input{figure19.eps_tex}}

\caption{Scheme of deformation of curve}
\label{Scheme of deformation of curve}
\end{figure}
From \eqref{Parametrise_1}, \eqref{Parametrise_2} one gets

\begin{align*}
d\accentset{\ast}{s}(s) &= \sqrt{d\tenq{y}\cdot d\tenq{y}} =  \sqrt{\bar{\tenq{y}}'(s)\cdot\bar{\tenq{y}}'(s)}ds\\
\Rightarrow \accentset{\ast}{s}'(s) &= \sqrt{\bar{\tenq{y}}'(s)\cdot\bar{\tenq{y}}'(s)}\\
\Rightarrow \accentset{\ast}{s}'(0) &= \sqrt{\bar{\tenq{y}}'(0)\cdot\bar{\tenq{y}}'(0)}\\
&\vc{=}{Chain}{Rule} \sqrt{{\left. {\frac{{\partial \hat{y_k}}}{{\partial {x_p}}}} \right|_{\tenq{x} = \mathring{\tenq{x}}}} {\left. {\frac{d \bar{x}_p(s)}{ds}} \right|_{s=0}} 
{\left. {\frac{{\partial \hat{y_k}}}{{\partial {x_g}}}} \right|_{\tenq{x} = \mathring{\tenq{x}}}} {\left. {\frac{d \bar{x}_g(s)}{ds}} \right|_{s=0}}}\\
&= \sqrt{\hat{y}_{k,p}(\mathring{\tenq{x}})\bar{x}'_p(0)\hat{y}_{k,g}(\mathring{\tenq{x}})\bar{x}'_g(0)}\\
&= \sqrt{\tenq{F}(\mathring{\tenq{x}})\bar{\tenq{x}}'(0)\cdot\tenq{F}(\mathring{\tenq{x}})\bar{\tenq{x}}'(0)}\\
&= \sqrt{\left\{\tenq{F}^T(\mathring{\tenq{x}})\tenq{F}(\mathring{\tenq{x}})\tenq{e}_i\right\}\cdot\tenq{e}_i}\text{ (no sum)}
\end{align*}
Recall $\bar{\tenq{x}}'(0)=\tenq{e}_i\text{ for }\Gamma_i$

\begin{align*}
\accentset{\ast}{s}'(0) &= \sqrt{\left\{\left(2\tenq{\phi}(\mathring{\tenq{x}}) + \tenq{1}\right)\tenq{e}_i\right\}\cdot\tenq{e}_i} = \sqrt{2\phi_{ii}(\mathring{\tenq{x}}) + 1 }\text{ (no sum)}\\
&=\footnotemark \lambda(\mathring{\tenq{x}}, \tenq{e}_i)
\end{align*}
\footnotetext{by result of homogeneous definition}
\end{proof}

\item
Let $\Gamma_i$, $\Gamma_j$ ($i\neq j$), ($i,j$ fixed) be two regular arcs both issuing from $\mathring{\tenq{x}}$ and having $\tenq{e}_i$, $\tenq{e}_j$ as respective tangent vectors at $\mathring{\tenq{x}}$. 

Let 
\[\accentset{\ast}{\Gamma}_i = \hat{\tenq{y}}(\Gamma_i), \accentset{\ast}{\Gamma}_j = \hat{\tenq{y}}(\Gamma_j)\]

\begin{figure}[H]
\centering
\resizebox{0.9\textwidth}{!}{\input{figure20.eps_tex}}
\caption{Scheme of deformation of curves, $i \neq j$}
\label{Scheme of deformation of curve, i does not equal to j}
\end{figure}

Call $\accentset{\ast}{\theta}(\mathring{\tenq{x}}, \tenq{e}_i, \tenq{e}_j) \in [0,\pi]$ the angle between the tangent vectors of $\accentset{\ast}{\Gamma}_1$ and $\accentset{\ast}{\Gamma}_2$ at $\mathring{\tenq{y}} = \hat{\tenq{y}}(\mathring{\tenq{x}})$, and set
\[\mathring{\mu}_{ij}\equiv{\mu}_{ij}(\mathring{\tenq{x}}) = \frac{\pi}{2} - \accentset{\ast}{\theta}(\mathring{\tenq{x}}, \tenq{e}_i, \tenq{e}_j)\phantom{-}(i \neq j)\]
Then 
\begin{equation}
\sin\mathring{\mu}_{ij} = \frac{2\mathring{\phi}_{ij}}{\sqrt{1+2\mathring{\phi_{ii}}}\sqrt{1+2\mathring{\phi_{jj}}}}\phantom{-}(i \neq j)\text{ (no sum)} \tag{$\mathsection\mathsection$} \label{sin2}
\end{equation}

Proof equivalent to part(a) [take as exercise]

Note the $\eqref{sin2}$ coincides with the formula for $\sin\mu_{ij}$ appropriate to the approximately local homogeneous deformation. 
\end{enumerate}

\begin{note}

We also had proved that any deformation is locally homogeneous. Any deformation in general is \emph{not} in the form of 
\[\tenq{y}=\hat{\tenq{y}}(\tenq{x})=\tenq{F}\tenq{x}+\tenq{c}\]
or 
\[\textit{grad }\hat{\tenq{y}}(\tenq{x}) = \nabla \hat{\tenq{y}} = \accentset{\ast}{\tenq{F}}(\tenq{x})\]
i.e. $\nabla \hat{\tenq{y}}$ is function of $\tenq{x}$

\end{note}

\section{Infinitesimal deformation}

Consider

\[\hat{\tenq{y}}(\tenq{x}) = \tenq{F}\tenq{x}+\tenq{c} = \tenq{x} + \left(\tenq{F}-\tenq{1}\right)\tenq{x}+\tenq{c} = \tenq{u}(\tenq{x})+\tenq{x}\]
where
\[\tenq{u}(\tenq{x}) = \left(\tenq{F}-\tenq{1}\right)\tenq{x}+\tenq{c}\]
\[\nabla \tenq{u} = \tenq{F} - \tenq{1} = \footnotemark \text{ constant}\]
\footnotetext{by homogeneous definition}

\begin{figure}[H]
\centering
\resizebox{0.6\textwidth}{!}{\input{figure21.eps_tex}}

\caption{Scheme of infinitesimal deformation}
\label{Scheme of infinitesimal deformation}
\end{figure}


define 

\[\tenq{\varepsilon} = \frac{1}{2}\left(\nabla \tenq{u} + \left(\nabla \tenq{u}\right)^T\right) = \textit{sym }\nabla \tenq{u}\]
or
\[\varepsilon_{ij} = \frac{1}{2}\left(u_{i,j} + u_{j,i}\right)\]

\[\tenq{\omega} = \frac{1}{2}\left(\nabla \tenq{u} - \left(\nabla \tenq{u}\right)^T\right) = \textit{skew }\nabla \tenq{u}\]
or
\[\omega_{ij} = \frac{1}{2}\left(u_{i,j} - u_{j,i}\right)\]


\[\theta = \tr \tenq{\varepsilon} = \varepsilon_{ii} = u_{i,i}\cdots\text{infinitesimal volume dilatation}\]

Assume we have deformation such that 
\[\abs{\nabla\tenq{u}}\ll 1\text{ or }\sqrt{u_{i,j}u_{i,j}}\ll 1\]
i.e., infinitesimal deformation

\begin{align*}
\tenq{E} = \frac{1}{2}\left(\tenq{F}^T\tenq{F}-\tenq{1}\right) 
&= \frac{1}{2}\left[\left(1+\nabla\tenq{u}\right)\left(1+\left(\nabla\tenq{u}\right)^T\right) - \tenq{1}\right]\\
&= \frac{1}{2}\left[\nabla\tenq{u}+\left(\nabla\tenq{u}\right)^T + \cancelto{\text{neglected}}{\left(\nabla\tenq{u}\right)^T\left(\nabla\tenq{u}\right)}\right]\\
&\approx \frac{1}{2}\left[\nabla\tenq{u}+\left(\nabla\tenq{u}\right)^T \right] = \tenq{\varepsilon}
\end{align*}

\begin{theorem}\mbox{}\\

\[\epsilon(\tenq{e}_i) = \left[1+2E_{ii}\right]^{1/2}-1 \text{ (no sum)}\]

If $\abs{\nabla\tenq{u}}\ll 1$, then 
\[\epsilon(\tenq{e}_i)\sim E_{ii} \sim \varepsilon_{ii} \text{ (no sum)}\]

Also 

\[\sin\mu_{ij} = \frac{2E_{ij}}{\sqrt{1+2E_{ii}}\sqrt{1+2E_{jj}}}, i\neq j \text{ (no sum)}\]

then 

\[\mu_{ij}\sim 2E_{ij} \sim 2\varepsilon_{ij} , i\neq j\]
and the volume dilatation

\begin{align*}
\Theta 
&= \frac{\volume(R_\ast) - \volume(R)}{\volume(R)} = J\footnotemark -1 = \det(\tenq{1}+\nabla\tenq{u}) - 1\\
&=\begin{vmatrix}
1+u_{1,1} & u_{1,2} & u_{1,3}\\
u_{2,1} & 1+u_{2,2} & u_{2,3}\\
u_{3,1} & u_{3,2} & 1+u_{3,3}
\end{vmatrix}  - 1\\
&\approx u_{1,1}+u_{2,2}+u_{3,3}=u_{i,i}=\tr \nabla\tenq{u} = \tr  \tenq{\varepsilon}
\end{align*}
\footnotetext{$J = \det\nabla\hat{\tenq{y}}=\det\tenq{F}$}

\end{theorem}

\begin{note}
For rigid deformation, $\tenq{y}=\tenq{F}\tenq{x}=\tenq{Q}\tenq{x}$
\[\tenq{E} = \frac{1}{2}\left(\tenq{F}^T\tenq{F}-\tenq{1}\right) = \frac{1}{2}\left(\tenq{Q}^T\tenq{Q}-\tenq{1}\right) = \tenq{0}\]

but 

\begin{align*}
\tenq{\varepsilon} 
&= \frac{1}{2}\left(\nabla\tenq{u}+\left(\nabla\tenq{u}\right)^T\right)\\
&= \frac{1}{2}\left(\tenq{Q}-\tenq{1}+\tenq{Q}^T-\tenq{1}\right)\\
&= \frac{1}{2}\left(\tenq{Q}+\tenq{Q}^T-2\tenq{1}\right)\neq\tenq{0}
\end{align*}
in general (except $\tenq{Q}\sim\tenq{1}$ or $\tenq{x}\sim\tenq{y}$)\\

Therefore, $\tenq{\varepsilon}$ is not good for rigid deformation. \footnote{In FEM, we always fixed the body because the basic formulas, i.e. shape functions $N$, are based on \emph{infinitesimal deformation.}}

\end{note}